
\begin{obeassessment}
\textbf{Kuis Bab 10}
1. Jelaskan peran fungsi Totient Euler dalam pembangkitan kunci RSA.
2. Bagaimana prosedur dekripsi RSA mengembalikan pesan asli menggunakan Teorema Euler?
3. Mengapa distribusi kunci publik rentan terhadap serangan \textit{Man-in-the-Middle}?
\end{obeassessment}
\begin{competencychecklist}
\begin{tabularx}{\linewidth}{|X|c|}
\hline
Indikator Kompetensi & Tercapai \\
\hline
Saya dapat membedakan kriptografi simetri dan asimetris & [ ] \\
\hline
Saya dapat menghitung pasangan kunci RSA untuk parameter kecil & [ ] \\
\hline
Saya dapat melakukan proses enkripsi dan dekripsi RSA & [ ] \\
\hline
Saya dapat menjelaskan kaitan pemfaktoran bilangan bulat dengan keamanan RSA & [ ] \\
\hline
Saya dapat menganalisis mekanisme kriptografi hibrida & [ ] \\
\hline
\end{tabularx}
\end{competencychecklist}
